\documentclass[preview]{standalone}

\usepackage{amsmath}
\usepackage{amssymb}
\usepackage{bettelini}
\usepackage{stellar}
\usepackage{definitions}

\begin{document}

\id{measuretheory-measurable-functions}
\genpage

\section{Measurable maps}

\includesnpt{measurable-space-definition}

\begin{snippetdefinition}{measurable-map-definition}{Measurable map}
    Let \((X,\Sigma_X)\) and \((Y,\Sigma_Y)\) be \measurablespace[measurable spaces].
    A map \(f\colon X\fromto Y\) is \emph{measurable} if
    \[
        f^{-1}(A)\in\Sigma_X
        \qquad\text{for every }A\in\Sigma_Y.
    \]
\end{snippetdefinition}

\begin{snippetproposition}{topological-target-measurable-map-criterion}{Topological target criterion}
    Let \((X,\Sigma)\) be a \measurablespace, let \((Y,\tau)\) be a
    \topologicalspace, and let \(f\colon X\fromto Y\). Then
    \(f\colon(X,\Sigma)\fromto(Y,\mathcal B(Y))\) is measurable if and only if
    \[
        f^{-1}(V)\in\Sigma
        \qquad\text{for every }V\in\tau.
    \]
\end{snippetproposition}

\begin{snippetproposition}{measurable-preimage-family-sigma-algebra}{Sets with measurable preimage form a \(\sigma\)-algebra}
    Let \((X,\Sigma)\) be a \measurablespace, let \(Y\) be a \set, and let
    \(f\colon X\fromto Y\). Then
    \[
        \mathcal F_f=\{A\subseteq Y\suchthat f^{-1}(A)\in\Sigma\}
    \]
    is a \(\sigma\)-algebra on \(Y\).
\end{snippetproposition}

\begin{snippetproof}{measurable-preimage-family-sigma-algebra-proof}{measurable-preimage-family-sigma-algebra}{Sets with measurable preimage form a \(\sigma\)-algebra}
    Since \(f^{-1}(Y)=X\), the family \(\mathcal F_f\) contains \(Y\).
    Preimages commute with complements and countable unions:
    \[
        f^{-1}(Y\difference A)=X\difference f^{-1}(A),
        \qquad
        f^{-1}\left(\bigcup_{n=1}^{\infty}A_n\right)
        =\bigcup_{n=1}^{\infty}f^{-1}(A_n).
    \]
    Thus \(\mathcal F_f\) is closed under both operations.
\end{snippetproof}

\begin{snippetproof}{topological-target-measurable-map-criterion-proof}{topological-target-measurable-map-criterion}{Topological target criterion}
    If \(f\) is measurable, the preimage of every open set is measurable because
    \(\tau\subseteq\mathcal B(Y)\). Conversely, the family of subsets of \(Y\)
    with measurable preimage is a \(\sigma\)-algebra. If it contains \(\tau\),
    it also contains \(\sigma(\tau)=\mathcal B(Y)\).
\end{snippetproof}

\begin{snippet}{measurable-functions-integration-motivation}
    Lebesgue integration is defined for measurable functions. Continuous maps
    between topological spaces are measurable with respect to their Borel
    \(\sigma\)-algebras, so this framework contains the continuous functions
    used in classical integration.
\end{snippet}

\section{Extended-real-valued functions}

\begin{snippettheorem}{extended-real-measurable-function-ray-criterion}{Measurability criterion for extended-real-valued functions}
    Let \((X,\Sigma)\) be a \measurablespace and let
    \(f\colon X\fromto\overline{\realnumbers}\), where
    \(\overline{\realnumbers}=[-\infty,+\infty]\) has its order topology.
    Then \(f\) is measurable if and only if
    \[
        \{x\in X\suchthat f(x)>a\}
        =f^{-1}((a,+\infty])\in\Sigma
        \qquad\text{for every }a\in\realnumbers.
    \]
\end{snippettheorem}

\begin{snippetproof}{extended-real-measurable-function-ray-criterion-proof}{extended-real-measurable-function-ray-criterion}{Measurability criterion for extended-real-valued functions}
    The forward implication follows because each ray \((a,+\infty]\) is open
    in the order topology. Conversely, the rays \((a,+\infty]\), with
    \(a\in\realnumbers\), generate the Borel \(\sigma\)-algebra of
    \(\overline{\realnumbers}\). The family of sets having measurable preimage
    under \(f\) is a \(\sigma\)-algebra and, by assumption, contains every such
    ray. It therefore contains \(\mathcal B(\overline{\realnumbers})\).
\end{snippetproof}

\begin{snippetcorollary}{extended-real-measurability-equivalent-rays}{Equivalent ray criteria}
    Let \((X,\Sigma)\) be a \measurablespace and let
    \(f\colon X\fromto\overline{\realnumbers}\). Each of the following
    conditions, required for every \(a\in\realnumbers\), is equivalent to the
    measurability of \(f\):
    \[
        \{f>a\}\in\Sigma,
        \qquad
        \{f\geq a\}\in\Sigma,
        \qquad
        \{f<a\}\in\Sigma,
        \qquad
        \{f\leq a\}\in\Sigma.
    \]
\end{snippetcorollary}

\begin{snippetcorollary}{measurable-function-fibers-measurable}{Fibres of a measurable function}
    Let \((X,\Sigma)\) be a \measurablespace and let
    \(f\colon X\fromto\overline{\realnumbers}\) be measurable. Then
    \[
        f^{-1}(\{a\})\in\Sigma
        \qquad\text{for every }a\in\overline{\realnumbers}.
    \]
\end{snippetcorollary}

\begin{snippetproof}{measurable-function-fibers-measurable-proof}{measurable-function-fibers-measurable}{Fibres of a measurable function}
    Every singleton is a Borel subset of \(\overline{\realnumbers}\), and the
    preimage of every Borel set under a measurable function is measurable.
\end{snippetproof}

\begin{snippetexample}{measurable-fibers-do-not-imply-measurability}{Measurable fibres do not imply measurability}
    Let \(A\subseteq\realnumbers\) be non-Borel and define
    \[
        f(x)=
        \begin{cases}
            e^x & x\in A,\\
            -e^x & x\notin A.
        \end{cases}
    \]
    Every fibre of \(f\) is either empty or a singleton, hence is Borel.
    Nevertheless,
    \[
        f^{-1}((0,+\infty))=A
    \]
    is not Borel, so \(f\) is not measurable.
\end{snippetexample}

\begin{snippetproposition}{indicator-function-measurability}{Measurability of an indicator function}
    Let \((X,\Sigma)\) be a \measurablespace and let \(A\subseteq X\). The
    indicator function \(\indicator_A\colon X\fromto\{0,1\}\) is measurable if
    and only if \(A\in\Sigma\).
\end{snippetproposition}

\section{Operations and limits}

\begin{snippetproposition}{measurable-map-composition}{Composition of measurable maps}
    Let \((X,\Sigma_X)\), \((Y,\Sigma_Y)\), and \((Z,\Sigma_Z)\) be measurable
    spaces. If \(f\colon X\fromto Y\) and \(g\colon Y\fromto Z\) are measurable,
    then \(g\circ f\colon X\fromto Z\) is measurable.
\end{snippetproposition}

\begin{snippetproof}{measurable-map-composition-proof}{measurable-map-composition}{Composition of measurable maps}
    For every \(A\in\Sigma_Z\),
    \[
        (g\circ f)^{-1}(A)=f^{-1}(g^{-1}(A)).
    \]
    Since \(g^{-1}(A)\in\Sigma_Y\), its preimage under \(f\) belongs to
    \(\Sigma_X\).
\end{snippetproof}

\begin{snippetcorollary}{continuous-postcomposition-preserves-measurability}{Continuous postcomposition preserves measurability}
    Let \((X,\Sigma)\) be a \measurablespace, let \(Y,Z\) be topological spaces,
    let \(g\colon X\fromto Y\) be measurable with respect to
    \(\mathcal B(Y)\), and let \(h\colon Y\fromto Z\) be continuous. Then
    \(h\circ g\colon X\fromto Z\) is measurable with respect to
    \(\mathcal B(Z)\).
\end{snippetcorollary}

\begin{snippetproposition}{arithmetic-operations-measurable-functions}{Arithmetic operations on measurable functions}
    Let \((X,\Sigma)\) be a \measurablespace and let
    \(f,g\colon X\fromto\realnumbers\) be measurable. Then \(f+g\), \(fg\),
    \(\max\{f,g\}\), and \(\min\{f,g\}\) are measurable.
\end{snippetproposition}

\begin{snippetproposition}{countable-extrema-measurable-functions}{Countable extrema and limits of measurable functions}
    Let \((X,\Sigma)\) be a \measurablespace and let
    \(f_n\colon X\fromto\overline{\realnumbers}\) be measurable for every
    \(n\in\naturalnumbers\). Then the pointwise functions
    \[
        \sup_n f_n,
        \qquad
        \inf_n f_n,
        \qquad
        \limsup_{n\to\infty}f_n,
        \qquad
        \liminf_{n\to\infty}f_n
    \]
    are measurable whenever the displayed extended-real operations are defined.
    In particular, every pointwise limit of measurable functions is measurable.
\end{snippetproposition}

\begin{snippetproof}{countable-extrema-measurable-functions-proof}{countable-extrema-measurable-functions}{Countable extrema and limits of measurable functions}
    For every \(a\in\realnumbers\),
    \[
        \left\{\sup_n f_n>a\right\}
        =\bigcup_{n=1}^{\infty}\{f_n>a\},
    \]
    so \(\sup_n f_n\) is measurable. Since negation is continuous,
    \(\inf_n f_n=-\sup_n(-f_n)\) is also measurable. Finally,
    \[
        \limsup_{n\to\infty}f_n
        =\inf_m\sup_{n\geq m}f_n,
        \qquad
        \liminf_{n\to\infty}f_n
        =\sup_m\inf_{n\geq m}f_n.
    \]
    These are countable infima and suprema of measurable functions. If
    \(f_n\to f\) pointwise, then \(f=\limsup f_n=\liminf f_n\).
\end{snippetproof}

\begin{snippetproposition}{absolute-value-measurable-function}{Absolute value of a measurable function}
    Let \((X,\Sigma)\) be a \measurablespace and let
    \(f\colon X\fromto\overline{\realnumbers}\) be measurable. Then \(\lvert f\rvert\)
    is measurable.
\end{snippetproposition}

\begin{snippetproof}{absolute-value-measurable-function-proof}{absolute-value-measurable-function}{Absolute value of a measurable function}
    The functions
    \[
        f^+=\max\{f,0\},
        \qquad
        f^-=-\min\{f,0\}
    \]
    are nonnegative and measurable. Since \(\lvert f\rvert=f^++f^-\), the
    absolute value is measurable.
\end{snippetproof}

\begin{snippetexample}{measurable-absolute-value-nonmeasurable-function}{The converse for absolute values fails}
    Let \((X,\Sigma)\) be a \measurablespace and let \(A\subseteq X\) be
    nonmeasurable. Define
    \[
        f=\indicator_A-\indicator_{X\difference A}.
    \]
    Then \(\lvert f\rvert=1\) is measurable, whereas \(f\) is not measurable.
\end{snippetexample}

\begin{snippetproposition}{complex-function-measurability-components}{Measurability of complex-valued functions}
    Let \((X,\Sigma)\) be a \measurablespace and let
    \(f\colon X\fromto\complexnumbers\). Then \(f\) is measurable if and only if
    \(\Re f\) and \(\Im f\) are measurable real-valued functions.
\end{snippetproposition}

\end{document}
